\documentclass[10pt,a4paper]{amsart}
\usepackage[margin=19mm]{geometry}
\usepackage{amsmath,amssymb,mathtools}
\usepackage[scr=rsfs,cal=euler]{mathalfa}
\usepackage{xfrac}
\usepackage[unicode,hidelinks,bookmarks=false]{hyperref}
\newcommand{\ie}{i.e.\ }
\newcommand{\cf}{cf.\ }
\newcommand{\resp}{respectively\ }
\newcommand{\Subsection}{\subsection}
\providecommand{\nobreakdash}{\nobreak\hbox{-}}
\setlength{\emergencystretch}{2em}
\multlinegap0pt
\usepackage{amssymb}
\usepackage{mathtools}
\usepackage[shortlabels]{enumitem}
\usepackage{bm}
\newtheorem{lem}{}
\usepackage{cleveref}
\crefname{lem}{}{s}
\title{Wasserstein stability and the nonsingular Borel lifting problem}
\author{Nachi Avraham-Re'em, George Peterzil}
\date{Source: arXiv:2606.03721v2 (CC BY 4.0)}
\begin{document}
\maketitle

\noindent\emph{Source and licence.} Wasserstein stability and the nonsingular Borel lifting problem. Version arXiv:2606.03721v2 (CC BY 4.0), distributed by arXiv under the Creative Commons Attribution 4.0 International licence, http://creativecommons.org/licenses/by/4.0/, as recorded in the arXiv licence metadata for this exact version and shown on the abstract page https://arxiv.org/abs/2606.03721v2. Full licence notice: \texttt{licenses/arxiv-2606.03721-CC-BY-4.0.txt}.\par\smallskip

\noindent\textit{Editorial scope.} Counted unit: the article's lem lem:martcont (source0.tex lines 513--523). The article's complete original proof of it is delivered with it (source0.tex lines 525--566, one \texttt{proof} environment of 42 lines). No mathematics is written, repaired or re-derived: the statement and the proof are the article's own lines, in the article's own notation, and no retained line is substituted or re-flowed.  Retained uncounted as the source-local context the proof needs: lines 509--511.  Nothing was omitted from any retained span, so the package carries no omission record.  No retained line contains a citation, so the wrapper carries no bibliography and no bibliography entry.  No retained line contains a figure environment, an image-inclusion command or drawing code, so no image is delivered and the package declares no asset dependency.  Editorial formatting only, in addition: an AMS \texttt{amsart} preamble that restates the article's own declarations and macros used by the retained lines, namely the environments \texttt{lem} on separate counters, together with the standard packages those lines need.  The retained environments print in this document as Lemma 1 in order of appearance, whereas the article prints the same items with the numbering of its own full document.  The complete native source of the article as distributed in the arXiv e-print for this exact version is bundled unchanged as \texttt{sources/201996/source0.tex}.
\begin{equation}\label{eq:Weil}
E_{i}f\left(g\right)=\int_{H_{i}}f\left(hg\right)dm_{H_{i}}\left(h\right)=\int_{H_{i}/H_{i-1}}\int_{H_{i-1}}f\left(\sigma_{i}\left(q\right)hg\right)dm_{H_{i-1}}\left(h\right)d\bar{m}_{i}\left(q\right)\quad m\text{-a.e. }g\in K.
\end{equation}

\begin{lem}\label{lem:martcont}
In the setup above, for every \(1\leq i\leq N\) the following hold.
\begin{itemize}
    \item For every \(1\)-Lipschitz function \(f:K\to\mathbb{R}\),
    \[\left\|E_{i}f-E_{i-1}f\right\|_{L^{\infty}\left(K\right)}\leq\rho_{i}.\]
    \item For every \(\phi\in L^{1}\left(K,m\right)\),
    \[\left\|E_{i}\phi-E_{i-1}\phi\right\|_{L^{1}\left(m\right)}\leq\int_{H_{i}/H_{i-1}}2\mathbf{V}_{m}\big(\sigma_{i}\left(q\right)^{-1}.\phi,\phi\big)d\bar{m}_{i}\left(q\right).\]
\end{itemize}
Consequently, for every density \(\phi\in\mathcal{D}\left(K,m\right)\),
\[\mathbf{W}_{d}\left(\phi m,m\right)\leq\sum\nolimits_{i=1}^{N}2\rho_{i}\int_{H_{i}/H_{i-1}}\mathbf{V}_{m}\big(\sigma_{i}\left(q\right)^{-1}.\phi,\phi\big)d\bar{m}_{i}\left(q\right).\]
\end{lem}

\begin{proof}[Proof of \cref{lem:martcont}]
Fix \(1\leq i\leq N\). Using Weil's integration formula \eqref{eq:Weil} and the right-invariance of \(d\), for every \(1\)-Lipschitz function \(f:K\to\mathbb{R}\) and \(m\)-a.e. \(g\in K\) we have
\begin{align*}
\left|E_{i}f\left(g\right)-E_{i-1}f\left(g\right)\right|
&\leq\int_{H_{i}/H_{i-1}}\int_{H_{i-1}}\left|f\left(\sigma_{i}\left(q\right)hg\right)-f\left(hg\right)\right|dm_{H_{i-1}}\left(h\right)d\bar{m}_{i}\left(q\right)\\
&\leq\int_{H_{i}/H_{i-1}}\int_{H_{i-1}}d\left(\sigma_{i}\left(q\right)hg,hg\right)dm_{H_{i-1}}\left(h\right)d\bar{m}_{i}\left(q\right)\\
&=\int_{H_{i}/H_{i-1}}d\left(e_{K},\sigma_{i}\left(q\right)\right)d\bar{m}_{i}\left(q\right)\\
&\leq\rho_{i}.
\end{align*}
This shows that \(\left\|E_{i}f-E_{i-1}f\right\|_{L^{\infty}\left(K\right)}\leq\rho_{i}\). Similarly, for every \(\phi\in L^{1}\left(K,m\right)\) we have
\begin{align*}
\left\|E_{i}\phi-E_{i-1}\phi\right\|_{L^{1}\left(m\right)}
&\leq\int_{H_{i}/H_{i-1}}\int_{H_{i-1}}\int_{K}\left|\phi\left(\sigma_{i}\left(q\right)hg\right)-\phi\left(hg\right)\right|dm\left(g\right)dm_{H_{i-1}}\left(h\right)d\bar{m}_{i}\left(q\right)\\
&=\int_{H_{i}/H_{i-1}}\big\|\sigma_{i}\left(q\right)^{-1}.\phi-\phi\big\|_{L^{1}\left(m\right)}d\bar{m}_{i}\left(q\right)=\int_{H_{i}/H_{i-1}}2\mathbf{V}_{m}\big(\sigma_{i}\left(q\right)^{-1}.\phi,\phi\big)d\bar{m}_{i}\left(q\right).
\end{align*}

Now let \(\phi\in\mathcal{D}\left(K,m\right)\), and fix an arbitrary \(1\)-Lipschitz function \(f:K\to\mathbb{R}\). Put
\[\Delta_{i}f:=E_{i-1}f-E_{i}f\quad\text{and}\quad\Delta_{i}\phi:=E_{i-1}\phi-E_{i}\phi.\]
Using \(E_{0}f=f\), \(E_{N}f=\int_{K}fdm\), \(E_{0}\phi=\phi\), and \(E_{N}\phi=1\), we have the telescoping identities
\[f-\int_{K}fdm=E_{0}f-E_{N}f=\sum\nolimits_{i=1}^{N}\Delta_{i}f,\quad\text{and}\quad\phi-1=E_{0}\phi-E_{N}\phi=\sum\nolimits_{i=1}^{N}\Delta_{i}\phi.\]
We now claim to have the martingale telescoping identity
\begin{equation}
\begin{aligned}
\int_{K}f\cdot\left(\phi-1\right)dm
&=\int_{K}\left(f-E_{N}f\right)\cdot\left(\phi-E_{N}\phi\right)dm\\
&=\sum\nolimits_{i,j=1}^{N}\int_{K}\Delta_{i}f\cdot\Delta_{j}\phi dm=\sum\nolimits_{i=1}^{N}\int_{K}\Delta_{i}f\cdot\Delta_{i}\phi dm.
\end{aligned}\label{eq:MTID}
\end{equation}
Let us justify this computation. The first equality follows because \(\int_{K}\left(\phi-1\right)dm=0\). The second equality follows by the above telescoping identities. For the third equality, we need the standard fact
\[\int_{K}\Delta_{i}f\cdot\Delta_{j}\phi dm=0\,\text{ for all }i\neq j;\]
let us consider the case \(i<j\) (the case \(j<i\) is analogous): by definition \(\Delta_{j}\phi\) is \(\mathcal{I}_{j-1}\)-measurable, and since \(i<j\) we have \(\mathcal{I}_{j-1}\subseteq\mathcal{I}_{i}\), and so \(\Delta_{j}\phi\) is also \(\mathcal{I}_{i}\)-measurable. Then by the law of total expectation
\[E_{i}\Delta_{i}f=E_{i}\left[E_{i-1}f-E_{i}f\right]=E_{i}f-E_{i}f=0,\]
and hence, using again the law of total expectation, we deduce that
\[\int_{K}\Delta_{i}f\cdot\Delta_{j}\phi dm=\int_{K}E_{i}\left[\Delta_{i}f\cdot\Delta_{j}\phi\right]dm=\int_{K}E_{i}\Delta_{i}f\cdot\Delta_{j}\phi dm=0.\]
This justifies the third equality, and so \eqref{eq:MTID} is established. We can now complete the proof:
\begin{align*}
\left|\mathbb{E}_{\phi m}\left[f\right]-\mathbb{E}_{m}\left[f\right]\right|	
&=\Big|\int_{K}f\cdot\left(\phi-1\right)dm\Big|\leq\sum\nolimits_{i=1}^{N}\Big|\int_{K}\Delta_{i}f\cdot\Delta_{i}\phi dm\Big|\\
&\leq\sum\nolimits_{i=1}^{N}\left\Vert\Delta_{i}f\right\Vert_{L^{\infty}\left(K\right)}\cdot\left\Vert\Delta_{i}\phi\right\Vert_{L^{1}\left(m\right)}\leq\sum\nolimits_{i=1}^{N}2\rho_{i}\int_{H_{i}/H_{i-1}}\mathbf{V}_{m}\big(\sigma_{i}\left(q\right)^{-1}.\phi,\phi\big)d\bar{m}_{i}\left(q\right),
\end{align*}
where the first inequality is by \eqref{eq:MTID}, and the third inequality is by the first two parts of the lemma. Finally, since \(f\) is arbitrary and the right hand-side is independent of \(f\), the desired inequality follows.
\end{proof}

\end{document}
